% !PRINTED 569
\weberpage{569}

%<<CONT
where $\nu$, like $\mu$, is an integer of the field $\Omega$ that is \gesperrt{not divisible by $\mu$.}

If this condition is fulfilled, then $g$ is a root of $P_\mu(x)$. But since $\wp(\mu u)$ becomes infinite in the second order in the $u$-plane for $u\equiv 0\ (\mathrm{modd}.\ \omega_1,\omega_2)$, $P_\mu(x)$ is divisible by $(x-g)^2$, unless
\[ \wp(u) - g \]
itself vanishes in the second order for $u=\nu:\mu$; on account of § 46, (18) this occurs when
\[ \wp'\left(\dfrac{\nu}{\mu}\right) = 0 \]
holds, that is, when $\nu:\mu$ is a half period, and thus
\weq{6}{g = e_1, e_2, e_3}
results.

\begin{quote}
{\small 1. \gesperrt{By (4), if $\omega_1=1$, $\omega_2=\omega$ is set, every integer of the field $\Omega$ is a period of $\wp(u)$. The converse does not hold; but since $A\omega$ is an integer in $\Omega$, every period must be converted into an integer by multiplication by $A$.}}
\end{quote}

From this it follows that two values
\[ g = \wp\left(\dfrac{\nu}{\mu}\right), \qquad g' = \wp\left(\dfrac{\nu'}{\mu}\right) \]
are equal to one another then and only then, when
\[ \nu \equiv \pm \nu' \pmod{\mu} \]
holds. For if $g=g'$, then either
\[ \dfrac{\nu+\nu'}{\mu} \quad \text{or} \quad \dfrac{\nu-\nu'}{\mu} \]
must be converted into an integer by multiplication by $A$. But $A$ is relatively prime to $\mu$, and consequently one of these two numbers must itself be an integer. The number of incongruent values of $\nu$, however, is equal to $N(\mu) = m$ (Bd.\ II, § 165).

If we therefore let $\nu$ run through a complete system of residues with respect to the modulus $\mu$, excluding zero, we obtain every value of $g$ twice, except when
\weq{7}{2\nu \equiv 0 \pmod{\mu},}
and in this case $\wp\left(\dfrac{\nu}{\mu}\right)$ is one of the values $e$. From this it follows:
%CONT>>
